\section{Introduction}
\label{sec:intro}

Memory is the substrate on which intelligent agents reason, adapt, and persist. Without mechanisms to store, retrieve, and evolve information across interactions, agents remain stateless responders---incapable of personalization, long-horizon planning, or continual improvement. The rapid proliferation of large language model (LLM)-based agents has made agent memory a central research area, with hundreds of papers published since 2023 alone.

\citet{DBLP:journals/corr/abs-2512-13564} provided the field's first comprehensive survey, ``Memory in the Age of AI Agents,'' synthesizing roughly 200 papers published through January 2026 into a unified $3\times 3 \times 3$ taxonomy. Their framework organizes agent memory along three axes: \textbf{Forms} (what carries memory: Token-level, Parametric, and Latent representations), \textbf{Functions} (why agents need memory: Factual, Experiential, and Working memory), and \textbf{Dynamics} (how memory evolves: Formation, Evolution, and Retrieval processes). This taxonomy has proven remarkably effective at classifying the design space and has been widely adopted by subsequent work.

However, the interval between January and June 2026 has seen an acceleration in both the volume and thematic diversity of agent memory research. Our curated catalog now spans 220 papers, including 21 works published since the original survey's cutoff. Critically, these new papers explore directions that the $3\times 3 \times 3$ taxonomy, being oriented around storage form, functional purpose, and lifecycle stage, was not designed to capture: temporal dynamics such as forgetting curves and sleep consolidation; multimodal memory that spans vision, language, audio, and embodiment; graph and structured knowledge representations; neurocognitively inspired architectures that explicitly model human memory systems; lifelong and continual learning mechanisms; new benchmarks and evaluation frameworks; and hybrid retrieval strategies that fuse multiple signals. These themes are cross-cutting---they appear across multiple form--function--dynamic cells---and demand additional analytical dimensions.

This paper makes three contributions:

\begin{enumerate}
  \item \textbf{A data-driven curation methodology for living surveys} (Section~\ref{sec:methodology}). We treat the paper list itself as a structured dataset, with a single \texttt{papers.yaml} file as the source of truth. Automated scripts validate entries, fetch metadata from arXiv and Semantic Scholar, discover new papers, and generate outputs including the README, a JSON export for an interactive web interface, and BibTeX references. A GitHub Actions workflow enforces validation on every contribution. This methodology is generalizable to any survey that must remain current.
  \item \textbf{An extended taxonomy with three new orthogonal dimensions} (Section~\ref{sec:taxonomy}). Building on the $3\times 3 \times 3$ foundation, we propose three additional axes: \textbf{Temporal Dynamics} (None, Decay-based, Consolidation-based, or Bi-temporal) capturing how memory evolves over time; \textbf{Modality} (Text-only, Multimodal-in, Multimodal-out, or Full-multimodal) capturing the sensory channels memory systems handle; and \textbf{Biological Inspiration} (None, Cognitive-metaphor, Neuro-inspired, or Brain-architecture) capturing the degree of fidelity to human memory. These dimensions are additive---they do not replace the original taxonomy but enrich it.
  \item \textbf{Gap analysis and research roadmap} (Section~\ref{sec:gap-analysis}). We identify persistently sparse cells in the taxonomy (Experiential/Latent, Working/Parametric), missing evaluation infrastructure (no unified benchmark, no temporal or multimodal benchmarks), and six open problems that define a concrete research agenda for the field.
\end{enumerate}

All materials supporting this paper are available at \url{github.com/tobias-weiss-ai-xr/agent-memory-research}, including the full paper list (\texttt{papers.yaml}), automation scripts, an interactive web interface at the GitHub Pages site, and the LaTeX source. We release the paper on Zenodo with a companion dataset DOI.

The remainder of this paper is organized as follows. Section~\ref{sec:background} recaps the original $3\times 3 \times 3$ taxonomy and its future directions. Section~\ref{sec:methodology} describes our data-driven curation pipeline. Section~\ref{sec:taxonomy} presents the extended taxonomy with seven cross-cutting themes and three new dimensions. Section~\ref{sec:landscape} surveys the 2026 landscape of new contributions. Section~\ref{sec:gap-analysis} identifies gaps and proposes a research roadmap. Section~\ref{sec:benchmarks} compares existing benchmarks. Section~\ref{sec:conclusion} concludes.
